%%%%%%%%%%%%%%%%%%%%%%%%%%%%%%%%%
%
%       EXERCÍCIO
%
%%%%%%%%%%%%%%%%%%%%%%%%%%%%%%%%%

\definevalorquestao

\begin{exercicioBanco}[\valorquestao]
A resistência (em toneladas) de vigas de concreto produzidas por uma empresa, comporta-se conforme a função de probabilidade abaixo:
%
\[
    \begin{array}{c|c|c|c|c|c}
         \text{Resistência} & 2 & 3 & 4 & 5 & 6 \\
         \hline
         p_{i} & 0,1 & 0,1 & 0,4 & 0,2 & 0,2
    \end{array}
\]
%
Admita que essas vigas são aprovadas para uso em construções se suportarem pelo menos \(3\) toneladas. De um grande lote fabricado pela empresa, escolhemos 15 vigas ao acaso. Qual é a probabilidade de:
%
\begin{enumerate}[(a)]
    \item Todas serem aptas para construções?
    \item No mínimo \(13\) serem aptas?
\end{enumerate}
%
\end{exercicioBanco}

